\section{Introduction}\label{sec:intro}

The Kochen--Specker (KS) theorem~\cite{KochenSpecker1967} shows that quantum observables cannot
be assigned definite values independently of the measurement context (for a review see
Ref.~\cite{Budroni2022}).  It is proved by exhibiting a finite set of rays in Hilbert space that
admits no $\{0,1\}$-valued assignment compatible with orthogonality.  Dimension three is the
smallest in which such sets exist.  (Smaller sets exist in higher dimensions, for example 20
vectors in dimension four~\cite{Kernaghan1994} and 36 in dimension
eight~\cite{KernaghanPeres1995}.)  In dimension three the smallest known set is the 31-ray
configuration of Conway and Kochen ($\CK$), reported by Peres~\cite{Peres1993}; the best lower
bound is 24, obtained for real vectors by Kirchweger, Peitl and
Szeider~\cite{KirchwegerPeitlSzeider2023} and for real and complex vectors by Li, Bright and
Ganesh~\cite{LiBrightGanesh2024}, improving the earlier 22 of Uijlen and
Westerbaan~\cite{UijlenWesterbaan2016}.  Trandafir and Cabello~\cite{TrandafirCabello2025} proved
that the orthogonality relations of $\CK$ have a unique realization up to unitary equivalence,
and that no set of 30 or fewer rays arises from one specified construction starting from the
Yu--Oh set; that 31 is the minimum without that assumption remains a conjecture.

\paragraph{Results.}
For a family of grids depending on parameters, the question of Section~\ref{sec:prolegomena} is
which parameter values give an uncolorable grid and how small a KS set the grid then contains.  For three families we
classify colorability completely and determine every minimum except at two classes of points of
the third family.

\begin{enumerate}[label=(\roman*)]
\item \emph{Grids with two coordinate values, $\{0,\pm1,\pm x\}$, $x\ne0$ complex}
(Theorem~\ref{thm:twosymbol}).
The grid is uncolorable only at isolated points: $x$ or $1/x$ one of $\pm2$, $\pm\sqrt2$,
$\pm\sqrt{-2}$, or $x$ a primitive third or sixth root of unity.  No KS set on such a grid has
fewer than 31 rays.
\item \emph{Real grids with three coordinate values, $\{0,\pm1,\pm x,\pm y\}$}
(Theorem~\ref{thm:threesymbol}).
After normalization to $1<x<y$, the grid is uncolorable on six curves, where two coordinate
values have ratio 2 or~$\sqrt2$, and at eight further points.  Again no KS set has fewer than
31 rays.
\item \emph{Conjugate-closed grids, with coordinate values $\{0,\pm1,\pm x,\pm x^{*}\}$, $x\ne0$ complex}
(Theorem~\ref{thm:conj}, Fig.~\ref{fig:circle}).  The grid is uncolorable exactly on two
circles, four lines through the origin with the origin removed, and four further points.  On
the circle $|x|^2=2$ one fixed list of 45 coordinate triples is uncolorable at every point and no shorter
fixed list is
(Theorem~\ref{thm:circle}); the circle $|x|^2=\tfrac12$ carries the 33-ray family of Gould and
Aravind.
\end{enumerate}

A fourth family has a discrete parameter: the grid of $0$ and the $n$th roots of unity is
uncolorable exactly when $6\mid n$ (Theorem~\ref{thm:6n}).  Section~\ref{sec:catalogue} adds
grids and completed grids that lie outside these families, each with exact counts and either a certified minimum or
a stated upper bound.

\paragraph{Standard of proof.}
Each of the three continuous-parameter classifications reduces, by an argument given in the
text, to finitely many cases that are decided in exact arithmetic by named scripts.  The
roots-of-unity theorem has a structural proof for every~$n$, with one finite computation at
$n=6$.  In addition the following are proved in the Lean
proof assistant and checked by its kernel (Section~\ref{sec:methods}): Theorem~\ref{thm:6n} for
$n\ge3$; the uncolorability of the 45 coordinate triples for every $x$ on the circle; and the minimum
KS subsets of six grids (31, 33, 33, 33, 43 and 47 rays).

\paragraph{Relation to earlier versions.}
Versions 1--9 of this paper proposed a classification by ``cancellation mechanisms'' and
described the alphabets with uncolorable pools (in the present terms, the coordinate values of
uncolorable grids) as isolated ``islands''.  That
proposal did not survive checking and was withdrawn in version~10;
Appendix~\ref{sec:corrections} lists the withdrawn claims.  Version~11 restated the results
that replaced it.  The present version adds the prolegomena and renames the objects as stated
there; the conclusions of the theorems, the arguments of the proofs and the numbers are
unchanged.  Since the uncolorable grids are not isolated,
we do not use the word ``island''.

\section{Definitions and methods}\label{sec:prelim}

\begin{definition}
A \emph{ray} in $\C^3$ is a one-dimensional subspace, written $[v]$ for a nonzero vector~$v$.
Rays $[v],[w]$ are \emph{orthogonal} if $\braket vw=\sum_k v_k^{*}w_k=0$.  A \emph{triad} (or
\emph{basis}) is a set of three mutually orthogonal rays.
\end{definition}

\begin{definition}
A \emph{coloring} of a set $S$ of rays is a map $f:S\to\{0,1\}$ such that (a)~no two orthogonal
rays of $S$ both receive~1, and (b)~every triad contained in $S$ has exactly one ray with value~1.
$S$ is \emph{KS-uncolorable}, or a \emph{KS set}, if it has no coloring.  Its \emph{orthogonality
graph} has the rays as vertices and the orthogonal pairs as edges; the triads are its triangles.
Both the edges and the triangles enter the definition, and an edge lying in no triangle of $S$
is still a constraint.  In the terms of Section~\ref{sec:rule}, a ray with $f=1$ is
\emph{colored} and a ray with $f=0$ is \emph{uncolored}.
\end{definition}

A coloring of $S$ restricts to a coloring of any subset, because every orthogonal pair and every
triad of the subset is one of~$S$:

\begin{lemma}[Restriction]\label{lem:restrict}
A set of rays that contains an uncolorable subset is uncolorable.
\end{lemma}

\begin{definition}
Let $\mathcal A\subset\C$ be a finite set containing~0, the \emph{coordinate values}.  A
\emph{coordinate triple} is a nonzero triple $v\in\mathcal A^3$, and the \emph{grid} of $\mathcal A$ is
the set of rays
\[
S(\mathcal A)=\{\,[v]: v\in\mathcal A^3\setminus\{0\}\,\}
\]
(Section~\ref{sec:grid}).  A grid ``with $k$ coordinate values'' has $k$ nonzero ones, a value and
its negative counted as one.  Where $\mathcal A$ depends on a parameter, a coordinate triple is an
expression in the parameter, evaluated at each value of it; the lists of coordinate triples
in Section~\ref{sec:conj} are lists of such expressions, and their colorability at a parameter
value is that of the rays they name there.  The \emph{completed grid} $\bar S(\mathcal A)$ is the smallest set of rays containing
$S(\mathcal A)$ and containing, with every orthogonal pair $[v],[w]$, the third ray of their
triad, $[(v\times w)^{*}]$ (the ordinary cross product for real vectors).  In all cases reported
here the completion was computed to closure and is finite; we do not know a general finiteness
criterion.
\end{definition}

In a completed grid every orthogonal pair lies in exactly one triad, so the numbers of pairs $E$
and triads $T$ satisfy $E=3T$.  This is a convenient check: the Eisenstein grid has
$E=174$, $T=22$, so it is not closed, contrary to a statement in version~9.

\paragraph{Three cautions on scope.}
(1)~Uncolorability is a property of a grid, not of a ring or a field.  Every ring that appears in
this paper contains $\Z$ and therefore contains the coordinates of $\CK$; the statement ``this
field supports a KS set'' is true of all of them and distinguishes nothing.  (2)~The size of a
generator is not an invariant of the grid: multiplying every coordinate triple by $x^{-1}$ shows
$S(\{0,\pm1,\pm x\})=S(\{0,\pm1,\pm x^{-1}\})$, and $S$ is also unchanged by $x\mapsto-x$.
(3)~A grid and its completed grid are different objects.  Their counts and minima need not
agree; inclusion gives only one-way implications (if the grid is uncolorable so is its
completed grid, and a KS subset of the grid bounds the minimum of the completed grid from
above).  Every count below says which is meant.

\subsection{Methods}\label{sec:methods}

\paragraph{Exact arithmetic.}  Coordinates are elements of a number field $\Q[t]/(m(t))$ stored
as rational coefficient vectors; rays are compared after division by the first nonzero
coordinate; a pair is orthogonal when the Hermitian product is the zero element.  The
cyclotomic grids use exponents modulo~$n$, and the proof of Theorem~\ref{thm:circle} uses
integer polynomial identities.  No count in Sections~\ref{sec:real}--\ref{sec:catalogue} depends
on a floating-point tolerance.  The only floating-point quantities in the paper are the flex
counts of Section~\ref{sec:rigidity}; they are labelled as numerical.  Floating-point
evaluation is used elsewhere only as a cross-check of exact computations.

\paragraph{Colorability.}  Colorability of a finite set is decided by a SAT
solver~\cite{AudemardSimon2018}: one variable per ray, a clause forbidding two 1s on each
orthogonal pair, and a clause requiring at least one 1 on each triad.

\paragraph{Certified minima.}  Let $P$ be an uncolorable grid or completed grid.  If $M\subseteq P$ is colorable
then, by Lemma~\ref{lem:restrict}, every uncolorable subset of $P$ contains a ray outside~$M$.
The certifier keeps a list of such sets $M$, each verified colorable by the SAT solver, and asks
a second solver for a set $R$ of fewer than $u$ rays meeting the complement of every listed~$M$,
where $u$ is the size of the best uncolorable subset found so far.  If $R$ is uncolorable, $u$
decreases; if $R$ is colorable it is enlarged to a colorable set $M\supseteq R$ and $M$ joins the
list.  When no such $R$ exists, no uncolorable subset has fewer than $u$ rays, and the minimum is
exactly~$u$.  The answer is correct whenever the SAT solver is.  The list of sets $M$ is in principle a
certificate that can be checked independently.  For six grids it is checked in Lean (below);
for Theorem~\ref{thm:circle} it is archived; the other minima are reproducible by re-running
the scripts.  Minima that the procedure did not finish within
its time budget are reported as upper bounds.

Two reductions speed up the runs for Theorems~\ref{thm:conj} and~\ref{thm:threesymbol}; both
keep the answer exact.  First, if a ray of an uncolorable set lies in no triad of that set,
deleting it leaves an uncolorable set (a coloring of the rest extends by giving the ray~0).  So a
minimum set lies in the \emph{core} of~$P$, obtained by deleting rays that lie in no triad
until none is left, and the search is run there.  Second, if
$g$ is a permutation of the rays preserving orthogonality and $M$ is colorable, so is $g(M)$, so
each colorable $M$ may be added to the list together with all its images under a group of such
permutations (signed permutations of the coordinates, and complex conjugation where the
coordinate values are closed under it; each is checked to preserve the graph before use).

The same enumeration, continued after the first
minimum set is found, lists \emph{all} minimum sets of~$P$.

\paragraph{Proofs checked in Lean.}  Three files in the repository are checked by the kernel of
Lean~4 (with the mathlib library) and use no axioms beyond the standard three.  Where a proof
rests on a finite set of rays, their coordinates are exact elements of a ring written as
integer tuples, orthogonality is recomputed inside Lean, and the certificates, which are
computed outside Lean, are checked there and not trusted.
\texttt{Cyclotomic6n.lean} proves Theorem~\ref{thm:6n} for $n\ge3$.
\texttt{NamedSets.lean} proves that the 43-ray set, the 47-ray set and, for every $x$ with
$|x|^2=2$, the 45 coordinate triples of Theorem~\ref{thm:circle} are uncolorable, each from a case-split
refutation.  \texttt{Minimality.lean} proves minima.  Its certificate for a grid with minimum
$u$ is the list of colorable sets $M$ described above, each with a coloring, together with a
branching proof that no set of fewer than $u$ rays meets the complement of every~$M$: a node
puts one ray in or out, and a leaf names complements that avoid the rays put in and are pairwise
disjoint once the rays put out are removed, so that a set meeting all of them needs a separate
ray for each.  A soundness theorem for this checker is proved once.  The result for each of six
grids, those with coordinate values $\{0,\pm1,\pm2\}$, $\{0,\pm1,\pm\sqrt2\}$,
$\{0,\pm1,\pm\sqrt{-2}\}$, $\{0,\pm1,\pm\omega,\pm\omega^2\}$ with $\omega=e^{2\pi i/3}$, $\{0,\pm1,\pm\alpha,\pm\alpha^{*}\}$ with
$\alpha=(1+\sqrt{-7})/2$, and
$\{0,\pm1,\pm\psi,\pm\psi^2\}$ with $\psi^3=\psi^2+1$, is a theorem with two halves: some $u$ of
the grid's vectors form an uncolorable set, and every set of fewer than $u$ of them has a
coloring.  Three limits should be stated.  First, in Lean each grid is an explicit list of
vectors; that the list consists of all rays whose coordinates are coordinate values is not
proved there.  Second, the Lean statements are about sets of vectors, not rays.  This loses nothing,
by an argument that is not itself formalized: a coloring of rays gives a coloring of any choice
of representative vectors and conversely, and each listed grid has one vector per ray (the
export script checks this), so the minimum for vectors is the minimum for rays.  Third, these
are minima within a grid; they
say nothing about KS sets outside it.

The procedure used in earlier versions (described there as ``OCUS'') was different and was not
sound: when a candidate was colorable, it required the next candidate to contain a ray that
completes a violated triad with \emph{two} rays of the current candidate, ignoring triads with
two or three rays outside it.  That can exclude genuine uncolorable subsets, and it did: it
reported 52 for the completed golden-ratio grid, which contains an uncolorable subset of~51.
None of the minima in this version comes from that code.
